\section*{Resumen ejecutivo}
\addcontentsline{toc}{section}{Resumen ejecutivo}

El presente trabajo continúa el análisis desarrollado en el Trabajo Práctico N.\textsuperscript{o} 1
sobre dos series horarias de infraestructura productiva de la empresa BodegaAI: la cantidad total
de solicitudes atendidas por el balanceador de carga (\emph{RequestCount}, serie ALB) y la cantidad de
solicitudes por destino del servicio de tienda (\emph{RequestCountPerTarget}, serie store-service). El
historial disponible se extiende de un mes, en el trabajo anterior, a aproximadamente tres meses
(2,063 observaciones horarias por serie, entre el 3 de julio y el 26 de septiembre de 2026, una vez
descartada la rampa de encendido del balanceador). Una tercera serie, a cargo de otro integrante del
grupo, se incorpora como espacio reservado en este informe.

El objetivo consiste en determinar si los modelos de aprendizaje automático (Machine Learning),
aprendizaje profundo (Deep Learning), modelos híbridos, herramientas de AutoML y modelos de fundación
superan, bajo un protocolo honesto de validación, el desempeño de los modelos clásicos utilizados en el
trabajo anterior (SARIMA). Se evaluaron 32 configuraciones por serie, agrupadas en diez familias:
referencias ingenuas, modelos estadísticos clásicos, la reestimación del SARIMA del Trabajo Práctico N.\textsuperscript{o} 1,
árboles de decisión con boosting y ensambles, redes neuronales profundas, Prophet y NeuralProphet, un
híbrido de descomposición estacional con boosting, dos plataformas de AutoML y un modelo de fundación
de series temporales (Chronos). La selección del modelo final se realizó exclusivamente por el error
absoluto medio (MAE) de validación entre los modelos elegibles, es decir, aquellos cuyo puntaje de
validación no estuvo contaminado por haber sido utilizado también para ajustar sus propios
hiperparámetros; el conjunto de prueba nunca se utilizó para seleccionar.

Para la serie ALB, el modelo seleccionado fue un AutoARIMA ajustado automáticamente (MASE de validación
de 0,475), apenas por debajo de la referencia estacional ingenua (0,489) y con una mejora relativa
considerable frente al SARIMA refit del trabajo anterior (0,923). Sin embargo, en el conjunto de prueba
el propio SARIMA del trabajo anterior resultó el más preciso de las 32 configuraciones evaluadas (MASE de
1,526, primer puesto), por delante de un Prophet ajustado con ventana libre de fuga de información (1,578) y del modelo
seleccionado (2,058, duodécimo puesto), lo que revela una consistencia baja y no significativa entre
validación y prueba para esta serie (correlación de Spearman $\rho=0,237$, $p=0,192$). Para la serie
store-service, el modelo seleccionado fue la propia referencia estacional ingenua con período 24 (MASE de
validación de 0,480): ningún modelo de Machine Learning, Deep Learning, híbrido o de fundación logró
superarla de manera no contaminada en validación. No obstante, en el conjunto de prueba el mejor
desempeño correspondió a NeuralProphet (MASE de 1,071, primer puesto), con un puntaje de validación
apenas mediocre (decimocuarto puesto); el modelo seleccionado quedó en el décimo puesto de prueba (MASE
de 1,443). La consistencia entre validación y prueba resultó, en este caso, moderada y significativa
($\rho=0,654$, $p<0,001$), aunque no evita divergencias puntuales como la de NeuralProphet. En síntesis,
el enfoque clásico del trabajo anterior no fue superado de manera consistente por los modelos más
recientes: los supera en la serie ALB en el conjunto de prueba, y en store-service es superado solo por
un modelo de la familia Prophet que la regla de selección honesta no habría elegido.

Las principales limitaciones del trabajo son la ventana única de validación de 48 horas, previa a una
intervención de despliegue que debió imputarse; el historial todavía acotado (menos de tres meses); la
presencia de modelos con puntaje de validación contaminado en el conjunto explorado, documentada pero no
utilizada para seleccionar; y una posible anomalía en las últimas horas registradas de la serie store-service, no
verificada como error de exportación de datos. Se concluye que, bajo un protocolo de validación estricto,
la superioridad de los modelos de aprendizaje automático y profundo sobre el enfoque clásico del trabajo
anterior no es uniforme: depende de la serie, del horizonte evaluado y del criterio de comparación
utilizado.
